\subsection{Basic Principle of Counting}

\begin{theorem}{Basic Principle of Counting}{basic-counting-principle}
  Suppose a first task can be performed in $m$ ways. Regardless of how the first task is performed, a second task can be performed in $n$ ways. Then both tasks can be performed in $mn$ ways.
\end{theorem}

The possible outcomes are ordered pairs:
\[
  \begin{array}{cccc}
    (1,1)  & (1,2)  & \cdots & (1,n)  \\
    (2,1)  & (2,2)  & \cdots & (2,n)  \\
    \vdots & \vdots & \ddots & \vdots \\
    (m,1)  & (m,2)  & \cdots & (m,n)
  \end{array}
\]

\begin{example}[Seating Students]
  How many arrangements are possible for $10$ students in a row of $10$ seats?
  \[
    \begin{array}{c|ccccc}
      \text{Seat number} & 1  & 2 & 3 & \cdots & 10 \\
      \text{Choices}     & 10 & 9 & 8 & \cdots & 1
    \end{array}
  \]
  By repeated application of \ref{thm:basic-counting-principle}, the total number of arrangements is
  \[
    10 \cdot 9 \cdot 8 \cdots 1 = 10!.
  \]
\end{example}

\subsection{Arrangements with Repeated Objects}

\begin{example}[Arranging MISSISSIPPI]
  How many distinguishable words can be formed using all the letters of \texttt{MISSISSIPPI}?

  Label the repeated letters:
  \[
    M,\ I_1,I_2,I_3,I_4,\ S_1,S_2,S_3,S_4,\ P_1,P_2.
  \]
  There are $11!$ arrangements of the labeled letters. Each word is counted $4!\,4!\,2!\,1!$ times, since permuting labels among identical letters leaves the word unchanged. Thus the number of distinguishable words is
  \[
    \frac{11!}{4!\,4!\,2!\,1!}.
  \]
\end{example}

\begin{theorem}{Arrangements with Repeated Objects}{repeated-object-arrangements}
  Suppose $n$ objects consist of $r$ distinct types, with $n_i$ identical objects of type $i$, where
  \[
    n_1+n_2+\cdots+n_r=n.
  \]
  The number of distinguishable arrangements is
  \[
    \frac{n!}{n_1!\,n_2!\cdots n_r!}
    = \binom{n}{n_1,n_2,\ldots,n_r}.
  \]
\end{theorem}

\subsection{Combinations}

\begin{example}[Choosing a Committee]
  How many $4$-person committees can be formed from a group of $10$ people?

  There are $10\cdot9\cdot8\cdot7$ ordered selections. Each committee is counted $4!$ times, so the number of committees is
  \[
    \frac{10\cdot9\cdot8\cdot7}{4!}
    = \frac{10!}{4!\,6!}
    = \binom{10}{6}
    = \binom{10}{4}.
  \]
\end{example}

To choose $k$ objects from $n$ distinct objects without replacement:
\[
  \begin{array}{c|cccc}
    \text{Selection number} & 1 & 2   & \cdots & k       \\
    \text{Choices}          & n & n-1 & \cdots & n-(k-1)
  \end{array}
\]
Each unordered selection occurs in $k!$ orders. Hence the number of selections is
\[
  \frac{n(n-1)\cdots(n-(k-1))}{k!}
  = \frac{n!}{k!\,(n-k)!}
  = \binom{n}{k}.
\]
